\chapter{Attention and Transformers}

\section{Introduction}

Recurrent models process sequences sequentially, which limits parallelization and makes long-range dependencies difficult to learn. Early sequence-to-sequence (seq2seq) models compress an entire input sequence into a single fixed-size vector, creating a severe information bottleneck.

Attention mechanisms and transformers overcome these limitations by allowing direct interactions between all elements of a sequence. Instead of compressing information into a single state, attention enables \emph{content-based retrieval} across the entire sequence.

\section{Why Recurrence Was Not Enough}

RNN-based seq2seq models map an input sequence to a single hidden vector:
\[
(x_1,\dots,x_T) \mapsto h_T.
\]

This creates:
\begin{itemize}
    \item a bottleneck limiting information capacity,
    \item difficulty in modeling long-range dependencies,
    \item sequential computation preventing parallelism.
\end{itemize}

Attention addresses these issues by allowing each output to access all input positions directly.

\section{Attention as Adaptive Weighting}

Attention computes a weighted combination of values based on similarity.

\subsection{Attention Formula}

Given:
\begin{itemize}
    \item query $q$,
    \item keys $k_i$,
    \item values $v_i$,
\end{itemize}
attention is:
\[
\mathrm{Attn}(q, K, V)
=
\sum_i \alpha_i v_i,
\quad
\alpha_i = \frac{\exp(q^\top k_i)}{\sum_j \exp(q^\top k_j)}.
\]

\subsection{Interpretation}

\begin{itemize}
    \item $q^\top k_i$ measures similarity,
    \item $\alpha_i$ selects relevant information,
    \item output is a weighted aggregation.
\end{itemize}

Thus attention is a form of content-based retrieval.

\section{Self-Attention and Token Interaction}

In self-attention, queries, keys, and values all come from the same sequence.

\subsection{Matrix Form}

Let $X \in \mathbb{R}^{n \times d}$. Define:
\[
Q = X W_Q,\quad K = X W_K,\quad V = X W_V.
\]

Then:
\[
\mathrm{Attn}(X)
=
\mathrm{softmax}\!\left(\frac{QK^\top}{\sqrt{d}}\right)V.
\]

\subsection{Interpretation}

\begin{itemize}
    \item $QK^\top$ computes all pairwise interactions,
    \item softmax normalizes attention weights,
    \item each token aggregates information from all others.
\end{itemize}

Thus self-attention enables direct token-to-token communication.

\section{Positional Encoding and Order Information}

Self-attention is permutation-invariant: it does not encode order by itself.

To incorporate sequence order, positional encodings are added:
\[
X \leftarrow X + P,
\]
where $P$ encodes position information.

\subsection{Types}

\begin{itemize}
    \item fixed sinusoidal encodings,
    \item learned embeddings.
\end{itemize}

This allows the model to distinguish positions in the sequence.

\section{Multi-Head Attention}

Instead of a single attention operation, transformers use multiple heads:
\[
\mathrm{MultiHead}(X) = \mathrm{concat}(\mathrm{head}_1,\dots,\mathrm{head}_h) W_O.
\]

Each head computes:
\[
\mathrm{head}_i = \mathrm{Attn}(X W_Q^{(i)}, X W_K^{(i)}, X W_V^{(i)}).
\]

\subsection{Interpretation}

\begin{itemize}
    \item different heads capture different relationships,
    \item the model attends to multiple patterns simultaneously.
\end{itemize}

\section{Transformer Blocks}

A transformer layer consists of:
\begin{itemize}
    \item multi-head self-attention,
    \item feedforward (MLP) layer,
    \item normalization,
    \item residual connections.
\end{itemize}

\subsection{Structure}

\[
X \to X + \mathrm{Attn}(X) \to \text{Norm} \to X + \mathrm{MLP}(X) \to \text{Norm}.
\]

This combines attention with nonlinear transformation and stabilization.

\section{Encoder, Decoder, and Causal Masking}

\subsection{Encoder}

Processes the entire sequence with full self-attention.

\subsection{Decoder}

Generates outputs sequentially, using \emph{causal masking} to prevent access to future tokens.

\begin{definition}[Causal Masking]
In attention weights, enforce:
\[
\alpha_{ij} = 0 \quad \text{for } j > i.
\]
\end{definition}

This ensures autoregressive behavior.

\subsection{Encoder--Decoder Attention}

The decoder attends both to:
\begin{itemize}
    \item its own past outputs,
    \item the encoder outputs.
\end{itemize}

\section{From Similarity to Aggregation (Worked Derivation)}

Starting from similarity scores:
\[
s_i = q^\top k_i,
\]
we normalize:
\[
\alpha_i = \frac{e^{s_i}}{\sum_j e^{s_j}},
\]
and compute:
\[
\mathrm{output} = \sum_i \alpha_i v_i.
\]

Thus attention is:
\[
\text{similarity} \to \text{weights} \to \text{aggregation}.
\]

\section{Scaling and Limitations}

\subsection{Advantages}

\begin{itemize}
    \item parallel computation,
    \item direct modeling of long-range dependencies,
    \item flexible interaction patterns.
\end{itemize}

\subsection{Limitations}

\begin{itemize}
    \item quadratic cost in sequence length,
    \item large memory usage,
    \item sensitivity to data scale.
\end{itemize}

These motivate efficient attention variants.

\section{A Unifying Perspective}

Attention replaces recurrence with direct interaction. Instead of passing information through a chain of states, it allows every element to access every other element in one step.

Transformers combine:
\begin{itemize}
    \item attention for interaction,
    \item MLPs for nonlinear transformation,
    \item normalization and residuals for stability.
\end{itemize}

This architecture has become a dominant paradigm in modern machine learning.

\section{Summary}

In this chapter we developed attention mechanisms and transformer architectures. We motivated attention from the limitations of recurrence, defined attention as content-based retrieval, and derived self-attention in matrix form.

We introduced positional encoding, multi-head attention, and transformer blocks, and explained encoder--decoder structures and causal masking. Transformers provide a powerful and flexible framework for sequence modeling.

\chapter*{Part II Exercises: Core Supervised Learning Models and Deep Learning}
\addcontentsline{toc}{chapter}{Part II Exercises}

\section*{Exercises}

\begin{enumerate}
    \item Compare linear models, kernel methods, and neural networks in terms of:
    \begin{enumerate}
        \item representation,
        \item inductive bias,
        \item scalability.
    \end{enumerate}

    \item Show why XOR cannot be solved by a linear classifier but can be solved by a two-layer network.

    \item Write a two-layer neural network explicitly as a composition of functions.

    \item Count the number of parameters in:
    \begin{enumerate}
        \item a fully connected layer,
        \item a convolutional layer.
    \end{enumerate}

    \item Derive backpropagation updates for a simple two-layer network.

    \item Explain why gradients vanish in deep networks.

    \item Compare ReLU and sigmoid activations in terms of gradient behavior.

    \item Explain why batch normalization improves optimization.

    \item Compare dropout and weight decay conceptually.

    \item Show how a residual block approximates an identity mapping.

    \item Prove translation equivariance of convolution in a simple 1D case.

    \item Compute the output size of a convolution layer given kernel size and stride.

    \item Unroll an RNN for $T$ steps and write it as a feedforward network.

    \item Derive the gradient of an RNN through time.

    \item Compare LSTM and GRU gating mechanisms.

    \item Compute a simple attention output for a toy example.

    \item Compare path length between two tokens in:
    \begin{enumerate}
        \item an RNN,
        \item a transformer.
    \end{enumerate}

    \item Explain why transformers handle long-range dependencies better than RNNs.

    \item Describe the role of positional encoding.

    \item Explain how multi-head attention increases model capacity.

    \item Design a simple architecture for:
    \begin{enumerate}
        \item image classification,
        \item sequence modeling,
        \item tabular data.
    \end{enumerate}
\end{enumerate}